\documentclass[9pt]{beamer}
%\documentclass[12pt,draft]{beamer} %Para dejar solo el contenido
\usetheme{Berlin}
%===================Temas predefinidos de colores=========================
%\usecolortheme{albatross}
%\usecolortheme{beaver}
%\usecolortheme{beetle}
%\usecolortheme{crane}
%===================Ajustes personalizados de colores ====================
\usepackage{xcolor}%Para ajustes sobre colores.  
\definecolor{Color1}{RGB}{0,200,0}
\definecolor{Color2}{RGB}{200,0,200}
\usepackage{comment}
\begin{comment}
%\setbeamercolor{title}{fg=Color1}
%Para ajustar el color del titulo, el primer plano se define a través de fg (foreground) y el fondo se define por bg (background)
\setbeamercolor{title}{fg=Color1, bg=Color2}% Ajuste del título.
%Ajuste de los títulos de los frames.
\setbeamercolor{frametitle}{bg=Color1, fg=Color2}
%Ajuste de la estructura restante del documento. 
\setbeamercolor{structure}{fg=yellow}
%Ajuste del color del texto de todas la presentación.
\setbeamercolor{normal text}{fg=blue,bg=cyan}
\end{comment}
%==========================Ajustes del tipo de letra===================== 
\usefonttheme{serif}
%Lista de opciones: default, serif, structurebold, structureitalicserif, structuresmallcapsserif, professionalfonts, serifbold, serifitalic
\usepackage[utf8]{inputenc}
\usepackage[spanish]{babel}
\usepackage{amsmath}
\usepackage{multicol}
\usepackage{amsfonts}
\usepackage{amssymb}
\usepackage{amsthm}%Paquete para teoremas
\usepackage{graphicx}
\usepackage{cancel}
\usepackage{tikz}
%================Personalizando Comandos==================================
\newcommand{\norma}[1]{\lVert #1 \rVert}
\newcommand{\norm}[1]{\left \lVert #1 \right \rVert}
\newcommand{\br}[1]{\left \lbrace #1 \right\rbrace}
\newcommand{\mb}[2]{\mathbb{#1}^#2}
%================Definiendo Teoremas======================================
\setbeamertemplate{theorems}[numbered]%Numerara teoremas
\newtheorem{Teo}{\textbf{Teorema}}
\newtheorem{Def}{\textbf{Definición}}[section]
\newtheorem{Prop}{\textbf{Proposición}}[section]
\newtheorem{Eje}{\textbf{Ejemplo}}[section]
%================Ajustes de video========================================
\usepackage{animate}
%ffmpeg -i file_example_MP4_480_1_5MG.mp4 Frames_%04.png
\usepackage{multimedia}
%================Titulo de la presentacion===============================
\title{Proyecto de Vinculación Euclides 4
}
\subtitle{Visita 4}
\author{\inst{1}Jose Mauricio Alvarenga Rodriguez}
%\setbeamercovered{transparent} 
%\setbeamertemplate{navigation symbols}{} 
%===============Ajuste del logo de la presentacion========================
%\logo{\includegraphics[scale=0.1]{Logo_UNAH_Gris}}
%\titlegraphic{\includegraphics[scale=0.1]{Logo_UNAH_Gris}}
\institute{\inst{1} Universidad Nacional Autónoma de Honduras} 
%\begin{comment}
\bfseries
\titlegraphic{ 
\begin{tikzpicture}[remember picture,overlay]
\node at (4.5,2){
    \includegraphics[scale=0.1]{Logo_UNAH_Gris}
};
\node at (-4.5,2){
    \includegraphics[scale=0.2]{ICVCLOGO}
};
\end{tikzpicture}
}
%\end{comment}
%\date{} 
%\subject{} 
\AtBeginSection[]{
\begin{frame}
\tableofcontents[currentsection]
\end{frame}
}
\begin{document}
\begin{frame}
\titlepage
\end{frame}
\section[Teoría Geometría]{Teoría sobre Geometría Básica}
\section[Ejemplos Geometría]{Ejemplos de Problemas de Geométría}
\begin{comment}
\begin{frame}{Ejemplos I}
\begin{multicols}{4}
\begin{enumerate}
\item[I.]\ \\
\scalebox{1}{
\begin{tikzpicture}
\draw (0,0)--(1,0)--(1,1)--(0,0);
\draw (0.8,0)--(0.8,0.2)--(1,0.2);
\draw (1.2,0.5) node[scale=0.7] {$4$}; 
\draw (0.5,-0.2) node[scale=0.7] {$3$}; 
\end{tikzpicture}}
\item[II.]\ \\
\scalebox{1}{
\begin{tikzpicture}
\draw (0,0)--(1.5,0)--(0.5,1)--(0,0);
\draw (1.3,0.5) node[scale=0.7] {$6$}; 
\draw (0.1,0.6) node[scale=0.7] {$5$}; 
\draw (0.3,0.2) node[scale=0.7] {$75^\circ$}; 
\end{tikzpicture}}
\item[III.]\ \\ 
\scalebox{1}{
\begin{tikzpicture}
\draw (0,0)--(1,-1)--(1,0)--(0,0);
\draw (0.8,0)--(0.8,-0.2)--(1,-0.2);
\draw (0.5,0.2) node[scale=0.7] {$6$}; 
\draw (1.1,-0.5) node[scale=0.7] {$8$}; 
\end{tikzpicture}}
\item[IV.]\ \\ 
\scalebox{1}{
\begin{tikzpicture}
\draw (0,0)--(0.5,1)--(1,0)--(0,0);
\draw (0.1,0.5) node[scale=0.7] {$8$}; 
\draw (0.9,0.5) node[scale=0.7] {$10$}; 
\draw (0.5,0.6) node[scale=0.7] {$75^\circ$}; 
\end{tikzpicture}}
\end{enumerate}
\end{multicols}
¿Cuáles de los siguientes triángulos anteriores son semejantes?
\begin{multicols}{4}
\begin{enumerate}
\item[A.] I y II
\item[B.] I y III
\item[C.] II y III
\item[D.] II y IV
\end{enumerate}
\end{multicols}
\onslide{2}{
\scalebox{1}{
\begin{tikzpicture}
\draw (0,0)--(1,0)--(1,1)--(0,0);
\draw (0.8,0)--(0.8,0.2)--(1,0.2);
\draw (1.2,0.5) node[scale=0.7] {$4$}; 
\draw (0.5,-0.2) node[scale=0.7] {$3$}; 
\draw (0,0)--(1,0)--(1,1)--(0,0);
\draw (0.8,0)--(0.8,0.2)--(1,0.2);
\draw (1.2,0.5) node[scale=0.7] {$4$}; 
\draw (0.5,-0.2) node[scale=0.7] {$3$}; 
\end{tikzpicture}}
}
\end{frame}
\end{comment}
\begin{frame}{Ejemplos I}
\begin{Eje}[Cálculo de perímetro de circunferencia]
¿Cuántos centímetros mide la circunferencia de un círculo cuyo radio mide 7 cm?
\end{Eje}
\begin{itemize}
\item<2-> \textcolor{blue}{Perímetro de Circunferencia=2$\times \pi\times $radio del círculo.}
\item<3-> \textcolor{blue}{Perímetro de Circunferencia=2$\times \pi\times 7cm=14\pi.$}
\end{itemize}
\onslide<4->{
\begin{Eje}[Área de círculo]
En la figura siguiente, el diámetro $AB$ mide 8cm. ¿Cuántos centímetros cuadrados mide el área del círculo?\\
\centering\begin{tikzpicture}
\draw (0,0) circle (0.5);
\draw (-0.5,0)--(0.5,0);
\draw (-0.7,0) node {$A$};
\draw (0.7,0) node {$B$};
\end{tikzpicture}
\end{Eje}}
\begin{itemize}
\item<5-> \textcolor{blue}{radio=diámetro/2=$\dfrac{8cm}{2}=4cm$.}
\item<6-> \textcolor{blue}{Área del cículo=$\pi\times radio^2=\pi(4cm)^2=16\pi cm^2.$}
\end{itemize}
\end{frame}
\begin{frame}{Ejemplos II}
\begin{Eje}[Cálculo del radio]
Si el área de un círculo es $100\pi u^2$, entonces ¿Cuál es su radio?
\end{Eje}
\textcolor{red}{
\begin{align*}
\visible<2->{100\pi u^2=&\pi r^2.}\\
\visible<3->{100\cancel{\pi} u^2=&\cancel{\pi} r^2.}\\
\visible<4->{100 u^2=& r^2.}\\
\visible<5->{\sqrt{100 u^2}=& \sqrt{r^2}.}\\
\visible<6->{10u=&r.}\\
\end{align*}}
\end{frame}
\begin{frame}{Ejemplos III}
\begin{Eje}[Construcción]
En la figura siguiente, ¿Es posible encontrar el radio?\\
\centering
\scalebox{0.7}{
\begin{tikzpicture}
\draw (0,0) circle (1);
\draw (-0.9,-0.43)--(-0.2,0.98);
\draw (-0.2,0.98)--(0.7,-0.71);
\draw (-0.8,0.2) node {$4$};
\draw (0.45,0.2) node {$6$};
\end{tikzpicture}}
\end{Eje}
\onslide<2->{\textcolor{red}{ Respuesta: NO.}}
\onslide<3->{
\begin{Eje}[Périmetro de un círculo]
En la figura siguiente, el diámetro $AB$ mide 8cm. ¿Cuántos centímetros mide el perímetro del círculo?\\
\centering\begin{tikzpicture}
\draw (0,0) circle (0.5);
\draw (-0.5,0)--(0.5,0);
\draw (-0.7,0) node {$A$};
\draw (0.7,0) node {$B$};
\end{tikzpicture}
\end{Eje}}
\textcolor{red}{
\onslide<4->{$P=2\pi r$}\onslide<5->{$=2\pi \dfrac{D}{2}$}\onslide<6->{$=2\pi \dfrac{8cm}{2}=8\pi cm.$}}
\end{frame}
\begin{frame}{Ejemplos IV}
\begin{Eje}[Périmetro de figura irregular]
Encuentre el perímetro de la siguiente figura.\\
\centering
\scalebox{0.4}{
\begin{tikzpicture}
\draw (0,0)--(2,0)--(5,4)--(5,5)--(2,5)--(0,4)--(0,0); 
\draw (1,-0.3) node[scale=1.5] {$2$};
\draw (-0.3,2) node[scale=1.5] {$4$};
\draw (1,4.8) node[scale=1.5] {$x$};
\draw (3.5,5.3) node[scale=1.5] {$3$};
\draw (5.3,4.5) node[scale=1.5] {$1$};
\draw (3.8,2) node[scale=1.5] {$5$};
\end{tikzpicture}
}
\end{Eje}
\begin{multicols}{2}
\onslide<2->{
\scalebox{0.4}{
\begin{tikzpicture}
\draw (0,0)--(2,0)--(5,4)--(5,5)--(2,5)--(0,4)--(0,0); 
\draw[dashed] (0,4)--(5,4);
\draw[dashed] (2,5)--(2,0);
\draw (1,-0.3) node[scale=1.5] {$2$};
\draw (1,3.7) node[scale=1.5] {$2$};
\draw (-0.3,2) node[scale=1.5] {$4$};
\draw (1,4.8) node[scale=1.5] {$x$};
\draw (3.5,5.3) node[scale=1.5] {$3$};
\draw (5.3,4.5) node[scale=1.5] {$1$};
\draw (2.3,4.5) node[scale=1.5] {$1$};
\draw (3.8,2) node[scale=1.5] {$5$};
\end{tikzpicture}
}
}\\
\textcolor{red}{
\begin{align*}
\visible<3->{x=&\sqrt{1^2+2^2}=\sqrt{5}.}\\
\visible<4->{P=&4+2+5+1+3+\sqrt{5}=15+\sqrt{5}.}
\end{align*}
}
\end{multicols}
\end{frame}
\begin{frame}{Ejemplos V}
\begin{Eje}[Cálculo de un lado de un triángulo]
Sea $x$ la medida de uno de los lados de un triángulo equilátero ¿Cuál es el valor de $x$ si el perímetro del triángulo es 54?
\end{Eje}
\begin{multicols}{2}
\textcolor{red}{
\visible<2->{
\scalebox{0.8}{
\begin{tikzpicture}
\draw (-1,0)--(1,0)--(0,1.9)--(-1,0); 
\draw (0,-0.3) node[scale=1.5] {$x$};
\draw (0.6,1.2) node[scale=1.5] {$x$};
\draw (-0.6,1.2) node[scale=1.5] {$x$};
\end{tikzpicture}
}}\\
\visible<3->{
$$x+x+x=3x=54\rightarrow x=\dfrac{54}{3}=18.$$}
}
\end{multicols}
\end{frame}
\section[Ejercicios Geometría]{Ejercicios  de Problemas de Geometría}
%\begin{frame}{Introducción al Cálculo}
%Las principales propiedades del ÁLGEBRA son:
%\begin{itemize}
%\item Propiedad de tricotomía. 
%\item  Propiedad del neutro para la suma.
%\end{itemize}
%\end{frame}
%\begin{frame}{Normas vectoriales}
%Una norma $\lVert x \rVert$ ($||x||$), $\norma{x}$.\\
%Algunos conjuntos conocidos $\mathbb{R}$, $\mathbb{C}$, $\mathbb{Q}$\\
%Espacios vectoriales: $\mb{R}{n}$.
%\end{frame}
%\begin{frame}{Referencias}
%\begin{thebibliography}{10}
%\bibitem{Goldbach1742}[Goldbach, 1742]
%Christian Goldbach.
%\newblock A problem we should try to solve before the ISPN ’43 deadline,
%\newblock \emph{Letter to Leonhard Euler}, 1742.
%\end{thebibliography}
%\end{frame}
\end{document}